% IoT Security seminar: 9 minutes of presentation + 1 minute discussion.
% Compile with pdfLaTeX. No external images or bibliography files required.
% Speaker notes: change "hide notes" to "show only notes" below.
\documentclass[aspectratio=169,12pt]{beamer}
\usepackage[utf8]{inputenc}
\usepackage[T1]{fontenc}
\usepackage{lmodern}
\usepackage{booktabs}
\usepackage{tabularx}
\usepackage{array}
\usepackage{microtype}
\usefonttheme{professionalfonts}
\setbeameroption{hide notes}

\definecolor{Navy}{HTML}{12324A}
\definecolor{Teal}{HTML}{007F86}
\definecolor{Muted}{HTML}{52616B}
\definecolor{Pale}{HTML}{EDF4F7}
\setbeamercolor{normal text}{fg=Navy,bg=white}
\setbeamercolor{structure}{fg=Teal}
\setbeamercolor{frametitle}{fg=Navy,bg=white}
\setbeamercolor{title}{fg=Navy}
\setbeamercolor{block title}{fg=Navy,bg=Pale}
\setbeamercolor{block body}{fg=Navy,bg=Pale}
\setbeamerfont{frametitle}{size=\Large,series=\bfseries}
\setbeamerfont{title}{size=\huge,series=\bfseries}
\setbeamerfont{subtitle}{size=\large}
\setbeamertemplate{navigation symbols}{}
\setbeamertemplate{itemize item}{\small\textbullet}
\setbeamersize{text margin left=9mm,text margin right=9mm}
\setbeamertemplate{footline}{%
  \hspace{9mm}\textcolor{Muted}{\scriptsize IoT Security}
  \hfill\textcolor{Muted}{\scriptsize\insertframenumber/\inserttotalframenumber}
  \hspace{9mm}\vspace{3mm}}
\setbeamertemplate{frametitle}{\vspace{3mm}\insertframetitle\par\vspace{1mm}}
\newcommand{\source}[1]{\par\vfill{\fontsize{6.5}{7.5}\selectfont\color{Muted}#1\par}}
\newcommand{\gap}{\par\vspace{2mm}}
\newcolumntype{Y}{>{\raggedright\arraybackslash}X}
\title{IoT Security:\\The Weakest Link}
\subtitle{Connected devices in homes, healthcare, and industry}
\author{Harrison Katantazi and Joshua Akhigbemen}
\date{}

\begin{document}

% 1. 15 seconds
\begin{frame}[plain]
\vspace{8mm}
{\color{Teal}\small TOPIC 9}\par
\gap
{\usebeamerfont{title}\inserttitle\par}
\gap
{\large\insertsubtitle\par}
\vspace{8mm}
{\small\insertauthor\par}
\gap
{\small Botnets, protocol security, and the EU Cyber Resilience Act}
\note{Opening, 15 seconds. IoT connects physical devices to digital services.
A low-cost device can become an entry point or part of a wider attack.
The discussion concerns how devices create risk across different settings.}
\end{frame}

% 2. 35 seconds
\begin{frame}{The central security question}
\large
How can one connected device create risk for a much larger system?
\gap
\normalsize
\begin{itemize}
\item A device depends on its firmware, network, and supporting services.
\item Weak authentication or missing updates can expose that system.
\item Consequences depend on what the device can access or control.
\end{itemize}
\gap
\textbf{Discussion focus:} Which protections belong in the product,
and which belong in its deployment?
\source{Sources: [S2] NISTIR 8259A; [S5] NIST SP 800-82 Rev. 3.}
\note{35 seconds. Define IoT as connected physical devices that sense,
communicate, or act. Explain that the camera, its management interface,
and its cloud service form one security problem. The title describes
a potential weak link. Risk depends on exposure and the surrounding controls.
Sources: https://doi.org/10.6028/NIST.IR.8259A and
https://doi.org/10.6028/NIST.SP.800-82r3.}
\end{frame}

% 3. 55 seconds
\begin{frame}{Different settings, different consequences}
\small
\renewcommand{\arraystretch}{1.25}
\begin{tabularx}{\linewidth}{@{}p{22mm}YY@{}}
\toprule
\textbf{Setting} & \textbf{Connected devices} & \textbf{Potential impact}\\
\midrule
Home & Cameras, door locks, smart appliances
& Loss of privacy or control\\
Healthcare & Patient monitors, connected treatment devices
& Exposure of health data or interruption of care\\
Industry & Sensors, gateways, connected controllers
& False process data, downtime, or safety consequences\\
\bottomrule
\end{tabularx}
\gap
\normalsize
Healthcare and industrial responses must consider physical safety
and service continuity.
\source{Illustrative examples. Sources: [S4] FDA; [S5] NIST SP 800-82 Rev. 3.}
\note{55 seconds. These are illustrative risks, not claims about particular
incidents. A camera affects privacy. A patient monitor supports clinical
decisions. An industrial sensor supports a physical process.
Patching or isolation can also interrupt essential services.
Sources: https://www.fda.gov/consumers/consumer-updates/medical-device-cybersecurity-what-you-need-know
and https://doi.org/10.6028/NIST.SP.800-82r3.}
\end{frame}

% 4. 45 seconds
\begin{frame}{Weaknesses across the device lifecycle}
\begin{description}
\item[Purchase] Unclear support periods or security capabilities.
\item[Installation] Default credentials and exposed management services.
\item[Operation] Excessive access and limited security visibility.
\item[Maintenance] Vulnerabilities persist when updates are unavailable
or difficult to deploy.
\end{description}
\gap
\textbf{Ownership matters:} Someone must track the device and its support status.
\source{Sources: [S1] Mirai study; [S2] NISTIR 8259A; [S5] OT security guidance.}
\note{45 seconds. Organise the weaknesses by lifecycle rather than treating
them as a random list. Ask who owns updates for a device installed several
years earlier. In healthcare and industry, maintenance needs coordination
with the people responsible for safe operation.
Sources: [S1], [S2], and [S5] on the references slide.}
\end{frame}

% 5. 65 seconds
\begin{frame}{Mirai: everyday devices in a botnet}
\begin{itemize}
\item In 2016, Mirai compromised embedded and IoT devices,
including cameras and routers.
\item It found exposed Telnet services and tried known credentials.
\item Compromised devices joined a remotely controlled botnet.
\item Attackers used their combined traffic for distributed
denial-of-service (DDoS) attacks.
\end{itemize}
\gap
{\large\bfseries A device owner and an attack victim can be different people.}
\source{[S1] Antonakakis et al., Understanding the Mirai Botnet, USENIX Security 2017.}
\note{65 seconds. A botnet is a collection of compromised devices controlled
by an attacker. DDoS means many sources overwhelm a service together.
The research documented Mirai attacks against targets including Dyn.
A household camera can participate while its owner notices little change.
This case illustrates the wider cost of weak device security.
Source: https://www.usenix.org/conference/usenixsecurity17/technical-sessions/presentation/antonakakis.}
\end{frame}

% 6. 65 seconds
\begin{frame}{Protocol security depends on configuration}
\small
MQTT uses a broker to pass messages between publishers and subscribers.
\gap
\renewcommand{\arraystretch}{1.35}
\begin{tabularx}{\linewidth}{@{}YYY@{}}
\toprule
\textbf{Weak deployment} & \textbf{Risk} & \textbf{Protection}\\
\midrule
Exposed Telnet login & Credential compromise & Disable unnecessary remote services\\
MQTT without TLS & Traffic disclosure or alteration & TLS and certificate validation\\
MQTT without client authentication & Unauthorised connections & Device authentication\\
Overly broad topic permissions & Unauthorised reads or commands & Topic-level authorisation\\
\bottomrule
\end{tabularx}
\source{Sources: [S1] Mirai study; [S3] OASIS MQTT 5.0, Section 5.}
\note{65 seconds. Use a temperature sensor as the publisher, a broker as
the intermediary, and a dashboard as the subscriber. MQTT supports secure
deployments when implementers configure the required protections.
Encryption protects traffic. Authentication establishes identity.
Authorisation controls which topics a device may read or publish.
Sources: [S1]; https://docs.oasis-open.org/mqtt/mqtt/v5.0/os/mqtt-v5.0-os.html, Section 5.}
\end{frame}

% 7. 55 seconds
\begin{frame}{Restricting what a compromised device can reach}
\small
\textbf{Illustrative factory policy}
\gap
\renewcommand{\arraystretch}{1.25}
\begin{tabularx}{\linewidth}{@{}p{29mm}Y>{\raggedright\arraybackslash}p{39mm}@{}}
\toprule
\textbf{Source} & \textbf{Destination} & \textbf{Policy}\\
\midrule
Network camera & Approved recording server & Required traffic only\\
Network camera & Production controller & Deny\\
Administrator & Camera management interface & Admin only\\
IoT segment & Other network segments & Deny unless justified\\
\bottomrule
\end{tabularx}
\gap
\normalsize
Segmentation needs enforced traffic rules and monitoring.
\source{Example policy based on [S5] NIST SP 800-82 Rev. 3.}
\note{55 seconds. The policy is an example, not a universal firewall template.
Give devices only the connections needed for their function.
A VLAN alone does not define which traffic can cross between networks.
Operational teams must verify that restrictions preserve the required process.
Source: https://doi.org/10.6028/NIST.SP.800-82r3.}
\end{frame}

% 8. 50 seconds
\begin{frame}{A device security baseline}
\small
NISTIR 8259A identifies six core capabilities.
\gap
\renewcommand{\arraystretch}{1.28}
\begin{tabularx}{\linewidth}{@{}p{49mm}Y@{}}
\toprule
\textbf{Capability} & \textbf{Procurement question}\\
\midrule
Device identification & Can we identify and inventory it?\\
Device configuration & Can authorised staff change settings securely?\\
Data protection & How does it protect stored and transmitted data?\\
Logical access to interfaces & Can we restrict interfaces and access?\\
Software update & Can it install authorised updates securely?\\
Cybersecurity state awareness & Can it report its security condition?\\
\bottomrule
\end{tabularx}
\source{[S2] NISTIR 8259A, IoT Device Cybersecurity Capability Core Baseline.}
\note{50 seconds. These capabilities turn vague claims about a secure
device into concrete questions for a supplier. Select requirements to suit
the environment. The baseline is a starting point, not a certification
or a guarantee that every deployment is safe.
Source: https://doi.org/10.6028/NIST.IR.8259A.}
\end{frame}

% 9. 55 seconds
\begin{frame}{The EU Cyber Resilience Act}
\textbf{Regulation (EU) 2024/2847}
\gap
\begin{itemize}
\item Covers many connected hardware and software products
supplied commercially on the EU market.
\item Requires product risk assessment and cybersecurity safeguards.
\item Requires vulnerability handling and security updates
during the support period.
\item Introduces conformity assessment and CE-marking requirements.
\end{itemize}
\gap
\small
\textbf{Healthcare distinction:} Products covered by the EU
Medical Device or In Vitro Diagnostic Medical Device Regulations
fall outside the CRA's scope.
\source{[S6] CRA, Article 2(2), Article 13, and Annex I; European Commission summary.}
\note{55 seconds. The CRA places obligations on manufacturers, with related
duties for importers and distributors. It concerns product security.
Do not imply that every connected medical device falls under this Act:
Article 2 excludes products covered by Regulations 2017/745 and 2017/746.
Generic connected products used in a hospital need their own scope assessment.
Sources: https://eur-lex.europa.eu/legal-content/EN/TXT/PDF/?uri=CELEX:32024R2847
and https://digital-strategy.ec.europa.eu/en/policies/cra-summary.}
\end{frame}

% 10. 40 seconds
\begin{frame}{CRA implementation and reporting}
\begin{description}
\item[10 December 2024] Regulation entered into force.
\item[11 September 2026] Manufacturer reporting obligations began.
\item[11 December 2027] Main product requirements apply.
\end{description}
\gap
\textbf{Reporting concerns actively exploited vulnerabilities
and severe incidents affecting product security.}
\gap
An early warning is due within \textbf{24 hours} of awareness,
followed by a notification within \textbf{72 hours}.
\source{[S6] CRA, Articles 14 and 71; European Commission reporting guidance.
Status checked 12 November 2026.}
\note{40 seconds. Distinguish entry into force from the dates obligations
apply. Follow-up reports also apply: for exploited vulnerabilities,
within 14 days after a corrective or mitigating measure is available;
for severe incidents, within one month after the 72-hour notification.
This slide gives the headline milestones.
Source: https://digital-strategy.ec.europa.eu/en/policies/cra-reporting.}
\end{frame}

% 11. 60 seconds plus discussion
\begin{frame}{Discussion: a factory camera vulnerability}
\textbf{Hypothetical case}
\gap
A factory discovers a critical camera vulnerability.
The fix requires a restart during production.
The camera network can currently reach production controllers.
\gap
\begin{enumerate}
\item What should the team restrict immediately?
\item How should it plan the update without disrupting safe operation?
\end{enumerate}
\gap
\textbf{Takeaway:} Product security and deployment controls
share responsibility for the outcome.
\source{Discussion scenario. Suggested reasoning draws on [S2], [S5], and [S6].}
\note{60 seconds of wrap-up, then allow roughly one minute for discussion.
Invite a response before offering an answer. A reasonable approach is
to assess process dependencies, restrict unnecessary camera access,
and arrange a tested update with operations staff.
Avoid presenting an immediate restart as universally safe.
Link the answer to procurement support and the network policy.
Sources: [S2], [S5], and [S6].}
\end{frame}

% 12. Reference slide, no scheduled speaking time
\begin{frame}{References}
\fontsize{8.2}{9.4}\selectfont
\begin{description}
\item[S1] Antonakakis et al. (2017).
\href{https://www.usenix.org/conference/usenixsecurity17/technical-sessions/presentation/antonakakis}
{\textit{Understanding the Mirai Botnet}}.
26th USENIX Security Symposium, pp.~1093--1110.

\item[S2] NIST (2020).
\href{https://doi.org/10.6028/NIST.IR.8259A}
{\textit{IoT Device Cybersecurity Capability Core Baseline}}.
NISTIR 8259A.

\item[S3] OASIS (2019).
\href{https://docs.oasis-open.org/mqtt/mqtt/v5.0/os/mqtt-v5.0-os.html}
{\textit{MQTT Version 5.0}}, Section~5.

\item[S4] U.S. FDA.
\href{https://www.fda.gov/consumers/consumer-updates/medical-device-cybersecurity-what-you-need-know}
{\textit{Medical Device Cybersecurity: What You Need to Know}}.

\item[S5] NIST (2023).
\href{https://doi.org/10.6028/NIST.SP.800-82r3}
{\textit{Guide to Operational Technology (OT) Security}}.
SP 800-82 Rev.~3.

\item[S6] European Parliament and Council (2024).
\href{https://eur-lex.europa.eu/legal-content/EN/TXT/PDF/?uri=CELEX:32024R2847}
{\textit{Regulation (EU) 2024/2847}}.
European Commission:
\href{https://digital-strategy.ec.europa.eu/en/policies/cra-summary}{CRA summary}
and \href{https://digital-strategy.ec.europa.eu/en/policies/cra-reporting}{reporting guidance}.
\end{description}
\note{Primary sources and official guidance underpin the presentation.
The discussion case and example traffic policy are illustrative.
The USENIX paper provides the research case study.
Use the source links for questions about evidence or legal scope.}
\end{frame}
\end{document}